% mode: ini
% Aimed at get_preamble_token, section 958: \span expanding the next token once, a run of \span,
% \span before something unexpandable, \span with a conditional; \tabskip in u and v parts, after
% \span, with \globaldefs positive and negative, so eq_define and geq_define; \tabskip after the
% alignment reads back what the group restored.
\input prelude
\def\m{a##b}
\def\mm{\m}
\def\t{\tabskip=2pt}
\def\x{x}
% Expanded once.
\setbox1\vbox{\halign{\span\m\cr y\cr}} \showbox1
% Expanded once only: \mm gives \m, which stays, so there is a missing #.
\setbox1\vbox{\halign{\span\mm\cr y\cr}} \showbox1
% Twice, by a run.
\setbox1\vbox{\halign{\span\span\mm\cr y\cr}} \showbox1
% \span then an unexpandable token, and a space, and \relax.
\setbox1\vbox{\halign{\span a#\span b&\span\relax#\span \cr y&z\cr}} \showbox1
% \span with a conditional: \iftrue goes, its \fi stays in the template unless expanded too.
\setbox1\vbox{\halign{\span\iftrue a#\span\fi\cr y\cr}} \showbox1
% \span before \tabskip: the loop falls out with \tabskip, and it is an assignment.
\tabskip=1pt
\setbox1\vbox{\halign{\span\tabskip=7pt#\tabskip=9pt&\t#\span\t\cr y&z\cr}} \showbox1
\showthe\tabskip
% The glue is what \tabskip is when each record is made.
\tabskip=1pt plus 1fil
\setbox1\vbox{\halign to 100pt{#\tabskip=2pt plus 2fil&#\tabskip=3pt minus 1pt&#\cr a&b&c\cr}} \showbox1
\showthe\tabskip
% An expansion that sets the tabskip in the middle of a template.
\tabskip=0pt
\setbox1\vbox{\halign{#\span\t&#\cr a&b\cr}} \showbox1
\showthe\tabskip
% \globaldefs: the assignment outlives the alignment.
\globaldefs=1
\setbox1\vbox{\halign{\tabskip=4pt#\cr y\cr}} \showbox1
\showthe\tabskip
\globaldefs=0 \tabskip=0pt
\globaldefs=-1
\setbox1\vbox{\begingroup\global\tabskip=5pt\halign{\global\tabskip=6pt#\cr y\cr}\endgroup} \showbox1
\showthe\tabskip
\globaldefs=0 \tabskip=0pt
% \global before \tabskip stays as a token of the template, and it then runs in each entry.
\setbox1\vbox{\halign{\global\tabskip=8pt#\cr y\cr}} \showbox1
\showthe\tabskip \tabskip=0pt
% \tabskip with a keyword scan that has to expand, and that backs up.
\def\pl{plus}
\setbox1\vbox{\halign{\tabskip=1pt\pl 1pt\relax#\cr y\cr}} \showbox1
\tabskip=0pt
% \span's expansion gives the \cr.
\def\e{\cr}
\setbox1\vbox{\halign{#\span\e y\cr}} \showbox1
\end
